% GDZ PPN236018647,native585middle;HKUST586.
\NumberedItem{114}
\noindent\hspace*{\parindent}Nov. 30. Felix fuit dies, quo multitudinem classium formar[um] binar[iarum] per triplicem methodum assignare largitum est nobis; puta:
\begin{enumerate}
\item per prod[uctum] infin[itum],
\item per aggregatum infinitum,
\item per aggregatum finitum cotangentium seu logarithm[orum] sinuum.
\end{enumerate}
\par\noindent\hfill Brun[ovici, 1800 Nov. 30.]
\par\medskip\begingroup\small
\noindent\hspace*{\parindent}Siehe Werke II, S. 285. Vergl. auch die folgende Nummer.
\par\noindent\hfill\textsc{Bachmann.}
\par\endgroup

\NumberedItem{115}
\noindent\hspace*{\parindent}Dec. 3. Methodum quartam ex omnibus simplicissimam deteximus pro\linebreak det[erminantibus] negativis ex sola multit[udine] numeror[um] \(\rho,\rho'\) etc. petitam, si \(Ax+\rho,Ax+\rho'\), etc. sunt formae lineares divisor[um] for[mae] \(\Box+D\).
\par\noindent\hfill Ibid. [Brunovici, 1800 Dec. 3.]
\par\medskip\begingroup\small
\noindent\hspace*{\parindent}Siehe Werke II, S. 286. Vergl. zu den Nummern 114 und 115 auch die oben S. 91 abgedruckten Aufzeichnungen und den Schluß des oben S. 235 abgedruckten Briefes von \textsc{Pfaff} vom 8. Dezember 1800, ferner Werke I, S. 476 die Bemerkung und ebenda S. 466 das Additamentum zu dem art. 306. X der \textit{Disquisitiones arithmeticae}, Werke I, S. 375, sowie endlich den Art. 27 des Aufsatzes »Über \textsc{Gauss}' zahlentheoretische Arbeiten« Werke X2, S. 66.
\par\noindent\hfill\textsc{Klein.\quad Bachmann.}
\par\endgroup
